\subsection{积分}

我们已经（II, p.~54, 定理 2）把区间 I 上的规则函数的原函数作为阶梯函数的原函数的一致极限而得到。这一程序可以稍有不同的方式表述：设 \(x_0\)、\(x\) 是 I 的满足 \(x_0 < x\) 的两个任意点；我们把区间 \([x_0, x]\) 的分割理解为任何区间序列 \([x_i, x_{i+1}]\)（以 \([x_0, x]\) 为并），其中 \((x_i)_{0 \leqslant i \leqslant n}\) 是 \([x_0, x]\) 的点的严格递增序列，满足 \(x_n = x\)。我们把相对于定义在 I 上的向量函数 \(\mathbf{f}\) 与由诸 \([x_i, x_{i+1}]\) 组成的分割的黎曼和理解为任何形如 \(\sum_{i=0}^{n-1} \mathbf{f}(t_i)(x_{i+1} - x_i)\) 的表达式，其中诸 \(t_i\) 属于 \([x_i, x_{i+1}]\)（对 \(0 \leqslant i \leqslant n-1\)）。于是有下列命题：

命题 5. 设 \(\mathbf{f}\) 是区间 I 上的规则函数，\(\mathbf{g}\) 是 \(\mathbf{f}\) 在 I 上的一个原函数，而 \([x_0, x]\) 是含于 I 的紧区间。对每个 \(\varepsilon > 0\) 存在数 \(\rho > 0\)，使得对 \([x_0, x]\) 分成长度为 \(\leqslant \rho\) 的区间的每个分割有

\[
\left\| \mathbf {g} (x) - \mathbf {g} (x _ {0}) - \sum_ {i = 0} ^ {n - 1} \mathbf {f} (t _ {i}) (x _ {i + 1} - x _ {i}) \right\| \leqslant \varepsilon\tag{1}
\]

对相对于此分割的每个黎曼和成立。

事实上，设 \(\mathbf{f}_1\) 是阶梯函数，使 \(\| \mathbf{f}(y) - \mathbf{f}_1(y)\| \leqslant \varepsilon\) 对每个 \(y\in [x_0,x]\) 成立；把 \(\mathbf{f}_1\) 在 I 上的一个原函数记作 \(\mathbf{g}_1\)，就有

\[
\left\| \mathbf {g} (x) - \mathbf {g} \left(x _ {0}\right) - \left(\mathbf {g} _ {1} (x) - \mathbf {g} _ {1} \left(x _ {0}\right)\right) \right\| \leqslant \varepsilon \left(x - x _ {0}\right)
\]

（依中值定理），而另一方面

\[
\left\| \sum_ {i = 0} ^ {n - 1} \mathbf {f} (t _ {i}) (x _ {i + 1} - x _ {i}) - \sum_ {i = 0} ^ {n - 1} \mathbf {f} _ {1} (t _ {i}) (x _ {i + 1} - x _ {i}) \right\| \leqslant \varepsilon (x - x _ {0}).
\]

因此只需就 \(\mathbf{f}\) 是阶梯函数的情形证明本命题。设 \((y_{k})_{1\leqslant k\leqslant m}\) 是 \(\mathbf{f}\) 在 \([x_0, x]\) 中的不连续点的严格递增有限序列。对 \([x_0, x]\) 分成长度为 \(\leqslant \rho\) 的区间的每个分割，点 \(y_{k}\) 中的每一个至多属于这些区间中的两个；因此至多只有 \(2m\) 个区间，在其上 \(\mathbf{f}\) 不取常值；但在这样一个区间 \([x_i, x_{i+1}]\) 上有

\[
\left\| \mathbf {g} \left(x _ {i + 1}\right) - \mathbf {g} \left(x _ {i}\right) - \mathbf {f} \left(t _ {i}\right) \left(x _ {i + 1} - x _ {i}\right) \right\| \leqslant 2 \mathrm{M} \left(x _ {i + 1} - x _ {i}\right)
\]

其中用 M 表示 \(\|f\|\) 在 \([x_{0}, x]\) 上的上确界；另一方面，当 f 在 \([x_{i}, x_{i+1}]\) 上取常值时有

\[
\mathbf {g} (x _ {i + 1}) - \mathbf {g} (x _ {i}) - \mathbf {f} (t _ {i}) (x _ {i + 1} - x _ {i}) = 0.
\]

于是可以看出，差 \(\left\|\mathbf{g}(x)-\mathbf{g}(x_{0})-\sum_{i=0}^{n-1}\mathbf{f}(t_{i})(x_{i+1}-x_{i})\right\|\) 不能超过 \(4Mm\rho\)；因此取 \(\rho\leqslant\varepsilon/4Mm\) 即得 (1)。

注记。1) 当 \(\mathbf{f}\) 连续时，命题 5 可以更简单地证明：由于 \(\mathbf{f}\) 在 \([x_0, x]\) 上一致连续，存在 \(\rho > 0\)，使得在每个长度为 \(\leqslant \rho\) 且含于 \([x_0, x]\) 的区间上 \(\mathbf{f}\) 的振幅为 \(\leqslant \frac{\varepsilon}{x - x_0}\)；对 \([x_0, x]\) 分成区间 \([x_i, x_{i+1}]\)（长度为 \(\leqslant \rho\)）的每个分割，以及对 \(t_i\) 在 \([x_i, x_{i+1}]\) 中的每个选取（对 \(0 \leqslant i \leqslant n-1\)），阶梯函数 \(\mathbf{f}_1\)（等于 \(\mathbf{f}(t_i)\) 于 \([x_i, x_{i+1}]\)（\(0 \leqslant i \leqslant n-1\)）上，且等于 \(\mathbf{f}(x)\) 于点 \(x\) 处）都满足：\(\| \mathbf{f}(y) - \mathbf{f}_1(y)\| \leqslant \frac{\varepsilon}{x - x_0}\) 在 \([x_0, x]\) 上成立；若 \(\mathbf{g}_1\) 是 \(\mathbf{f}_1\) 的一个原函数，则有 \(\mathbf{g}_1(x) - \mathbf{g}_1(x_0) = \sum_{i=0}^{n-1} \mathbf{f}(t_i)(x_{i+1} - x_i)\)，于是关系 (1) 由中值定理立得。

在本章的余下部分，我们将限于研究区间 I 上规则函数的原函数。对这样一个函数 \(\mathbf{f}\)（取值于 E 中）、一个原函数 \(\mathbf{g}\)（\(\mathbf{f}\) 的），以及 I 的两个任意点 \(x_0\)、\(x\)，E 的元素 \(\mathbf{g}(x) - \mathbf{g}(x_0)\)（显然，无论考虑哪个原函数 \(\mathbf{g}\)（\(\mathbf{f}\) 的），它都是同一个）称为函数 \(\mathbf{f}\) 从 \(x_0\) 到 \(x\) 的积分（或在紧区间 \([x_0, x]\) 上的积分），记作 \(\int_{x_0}^{x} \mathbf{f}(t) dt\) 或 \(\int_{x_0}^{x} \mathbf{f}\)。这个名称与记法源于 II, p.~57 的命题 5，它表明积分可以用黎曼和任意接近地逼近；更特别地，取 \([x_0, x]\) 分成相等区间的分割，可以写

\[
\frac {1}{x - x _ {0}} \int_ {x _ {0}} ^ {x} \mathbf {f} (t) d t = \lim _ {n \rightarrow \infty} \frac {1}{n} \sum_ {k = 0} ^ {n - 1} \mathbf {f} \left(x _ {0} + k \frac {x - x _ {0}}{n}\right).\tag{2}
\]

换言之，元素 \(\frac{1}{x-x_{0}}\int_{x_{0}}^{x}\mathbf{f}(t)dt\) 是 \([x_{0},x]\) 分成相等区间的分割的各区间的左端点处 f 的值的算术平均的极限；我们也把它称为函数 f 在区间 \([x_{0},x]\) 上的平均值（或中值）。

按定义，函数 \(x \mapsto \int_{x_{0}}^{x} \mathbf{f}(t) dt\) 无非就是在点 \(x_{0} \in I\) 处为零的 f 的原函数；我们也把它记作 \(\int_{x_{0}} \mathbf{f}(t) dt\) 或 \(\int_{x_{0}} f\)。

注记。2) 对定义在 I 上、取值于 E 中的任意函数 \(\mathbf{h}\)，元素 \(\mathbf{h}(x) - \mathbf{h}(x_0)\) 也写作 \(\left.\mathbf{h}(t)\right|_{x_0}^x\)；用这个记法可以看出，若 \(\mathbf{g}\) 是 I 上规则函数 \(\mathbf{f}\) 的任一原函数，则有

\[
\int_ {x _ {0}} ^ {x} \mathbf {f} (t) d t = \mathbf {g} (t) \left| _ {x _ {0}} ^ {x}. \right.\tag{3}
\]

\begin{enumerate}
\def\labelenumi{\arabic{enumi})}
\setcounter{enumi}{2}
\item
  表达式 \(\int_{x_{0}}^{x}\mathbf{f}(t)dt\) 与 \(\mathbf{g}(t)\bigg|_{x_{0}}^{x}\) 是缩写记号，代表其中出现字母 x、\(x_{0}\)、f、g 而不出现字母 t 的总体（参见~Set Theory, I, p.~15）；我们说在这些记号中 t 是一个“哑变量”；因此可以把 t 换成任何异于 x、\(x_{0}\)、f 与 g（以及可能进入这些记号出现的证明中的变量）的其他变量，而不改变如此得到的记号的意义（读者可以把这些记号与诸如 \(\sum_{i=1}^{n}x_{i}\) 或 \(\bigcup_{i}X_{i}\)（其中 i 同样是哑变量）之类的记号相比较）。
\item
  用黎曼和逼近积分与积分概念的历史起源之一，即面积的度量问题，密切相关。我们将在《积分》一册书中回到这一点，该册书专门讨论这个问题所引向的积分概念的种种推广；在这些推广中，待“积分”的函数未必定义在 R 的子集上；此外，即使处理的是（未必规则的）实变实函数 f，对它们可以定义积分 \(\int_{x_{0}}^{x}f(t)dt\)，函数 \(x\mapsto\int_{x_{0}}^{x}f(t)dt\) 也并不总是 f 的原函数，并且存在具有原函数却在我们所暗指的意义下不可“积分”的函数。
\end{enumerate}
% semantic-fragments: legacy-input-seam
\relax{}
